% EGA IV, Errata y Addenda, lista 3, n.º 53.
% Este componente se inserta después de 21.4.8 en el texto de lectura corregido
% y se repite literalmente bajo Err_IV, 53 en el apéndice de erratas impreso.

La proposición siguiente y sus corolarios nos fueron señalados por
M.~Raynaud.

\begin{proposition}[21.4.9]
\ifdefined\EGAIVSuppressErratumLabels\else\label{IV.21.4.9}\fi
Sean $X$ e $Y$ dos preesquemas localmente noetherianos, y sea
$f:X\to Y$ un morfismo fielmente plano. Supongamos además que $f$ es
cuasicompacto o que $f$ es abierto. Un divisor $D'$ sobre $X$ es de la
forma $f^*(D)$, con $D$ un divisor sobre $Y$, si y solo si se cumple la
condición siguiente: para todo $y\in Y$ tal que
$\operatorname{prof}(\sh{O}_{Y,y})\leq 1$, pongamos
\[
 Y_1=\operatorname{Spec}(\sh{O}_{Y,y}),\qquad
 X_1=X\times_Y Y_1,
 \qquad f_1=f_{(Y_1)}:X_1\to Y_1;
\]
entonces la imagen recíproca $D_1'$ de $D'$ sobre $X_1$ es de la forma
$f_1^*(D_1)$ para algún $D_1\in\operatorname{Div}(Y_1)$.
\end{proposition}

\oldpage[IV]{359}
\begin{proof}
La necesidad de la condición resulta inmediatamente de la transitividad de las
imágenes recíprocas de divisores, pues todos los morfismos considerados son planos
{\EGAReadLink{IV.21.4.4}{(21.4.4)}}. Para todo abierto $U$ de $Y$, denotemos por $f_U$ la
restricción $f^{-1}(U)\to U$ de $f$. Por {\EGAReadLink{IV.21.4.8}{(21.4.8)}}, la aplicación
\[
 D_U\longmapsto f_U^*(D_U)
\]
de $\operatorname{Div}(U)$ en
$\operatorname{Div}(f^{-1}(U))$ es inyectiva. Así, si $U_1,U_2$ son
dos abiertos de $Y$ y
\[
 D'|f^{-1}(U_1)=f_{U_1}^*(D_1),\qquad
 D'|f^{-1}(U_2)=f_{U_2}^*(D_2),
\]
entonces necesariamente
$D_1|U_1\cap U_2=D_2|U_1\cap U_2$. Se sigue inmediatamente que existe un
mayor abierto $U\subset Y$ para el cual
\[
 D'|f^{-1}(U)=f_U^*(D_U)
 \quad\text{con}\quad D_U\in\operatorname{Div}(U).
\]
Debemos demostrar que $U=Y$. Supongamos lo contrario, y sea $y$ un punto
maximal de $Y\setminus U$. Encontraremos un entorno abierto $V$ de
$y$ en $Y$ tal que
  $D'|f^{-1}(V)=f_V^*(D_V)$ para algún
$D_V\in\operatorname{Div}(V)$, contradiciendo la maximalidad de $U$.

\textup{I) Caso en que $f$ es cuasicompacto.}
Utilizando la cuasicompacidad de $f$ y {\EGAReadLink{IV.20.3.8}{(20.3.8)}}, por el método de
({\EGAReadLink{IV.8.1.2}{8.1.2}}, a), nos reducimos a demostrar ---con las notaciones del
enunciado, pero sin imponer inicialmente una cota a
$\operatorname{prof}(\sh{O}_{Y,y})$--- que $D_1'$ es de la forma
$f_1^*(D_1)$. Si $\operatorname{prof}(\sh{O}_{Y,y})\leq 1$, esto es
precisamente la hipótesis. Supongamos, pues, que
$\operatorname{prof}(\sh{O}_{Y,y})\geq 2$. Como $y$ es un punto
maximal de $Y\setminus U$, el abierto $W=Y_1\cap U$ de $Y_1$ es
$Y_1\setminus\{y\}$. Podemos, por tanto, limitarnos al caso
$Y=Y_1$ y $U=Y\setminus\{y\}$.

Por hipótesis existe un divisor $D_U\in\operatorname{Div}(U)$ tal
que $f_U^*(D_U)=D'|f^{-1}(U)$. El divisor $D_U$ es la clase de
equivalencia de un par $(\sh{L}_U,s)$, donde $\sh{L}_U$ es un
$\sh{O}_U$-módulo invertible y $s$ una sección meromorfa regular de
$\sh{L}_U$ sobre $U$ {\EGAReadLink{IV.21.2.11}{(21.2.11)}}. Por tanto, en virtud de
{\EGAReadLink{IV.21.4.1}{(21.4.1)}}, $D'|f^{-1}(U)$ es la clase de equivalencia de
$(\sh{L}'_U,s')$, donde
\[
 \sh{L}'_U=f_U^*(\sh{L}_U),\qquad s'=s\circ f_U.
\]
Sean $i:U\to Y$ y $j:f^{-1}(U)\to X$ las inmersiones canónicas,
y pongamos
\[
 \sh{L}=i_*(\sh{L}_U),\qquad
 \sh{L}'=j_*(\sh{L}'_U).
\]
Como $f$ es plano, $f^*(\sh{L})$ es canónicamente isomorfo a
$\sh{L}'$ {\EGAReadLink{IV.2.3.1}{(2.3.1)}}. Además, la planitud y la desigualdad
$\operatorname{prof}(\sh{O}_{Y,y})\geq 2$ implican
$\operatorname{prof}(\sh{O}_{X,x})\geq 2$ para todo
$x\in f^{-1}(y)$ {\EGAReadLink{IV.6.3.1}{(6.3.1)}}.

El divisor $D'$ es la clase de equivalencia de un par
$(\sh{L}'',s'')$, donde $\sh{L}''$ es un $\sh{O}_X$-módulo invertible
y $s''$ una sección meromorfa regular de $\sh{L}''$ sobre $X$.
Por consiguiente, $\sh{L}'_U$ es isomorfo a
$\sh{L}''|f^{-1}(U)$, y las observaciones precedentes, junto con
{\EGAReadLink{IV.5.9.8}{(5.9.8)}} y {\EGAReadLink{IV.5.10.5}{(5.10.5)}}, muestran que $\sh{L}'$ es un
$\sh{O}_X$-módulo invertible isomorfo a $\sh{L}''$. Como $f$ es
fielmente plano y cuasicompacto, se sigue que $\sh{L}$ es un
$\sh{O}_Y$-módulo invertible {\EGAReadLink{IV.2.5.2}{(2.5.2)}} y, por tanto, isomorfo a $\sh{O}_Y$, pues
$Y$ es un esquema local.

Además, como $\operatorname{prof}(\sh{O}_{Y,y})\geq 1$, el abierto
$U=Y\setminus\{y\}$ es esquemáticamente denso en $Y$
({\EGAReadLink{IV.20.2.13}{20.2.13}}, (iv)). La sección meromorfa regular $s$ de $\sh{L}_U$
sobre $U$ se prolonga, por tanto, a una sección meromorfa regular $s_0$ de
$\sh{L}$ sobre $Y$ {\EGAReadLink{IV.20.2.11}{(20.2.11)}}. La clase de equivalencia de
$(\sh{L},s_0)$ es así un divisor $D$ sobre $Y$. Evidentemente, $f^*(D)$ y
$D'$ inducen el mismo divisor sobre
$f^{-1}(U)=X\setminus f^{-1}(y)$; y, como
$\operatorname{prof}(\sh{O}_{X,x})\geq 2$ para todo
$x\in f^{-1}(y)$, {\EGAReadLink{IV.21.1.8}{(21.1.8)}} da $D'=f^*(D)$. Esto prueba el caso I.

\textup{II) Caso en que $f$ es abierto.}
Se trata de encontrar un entorno abierto $V$ de $y$ en $Y$ tal
que $D'|f^{-1}(V)=f_V^*(D_V)$ para algún
$D_V\in\operatorname{Div}(V)$, por lo que podemos suponer primero que $Y$ es afín.
Como $f$ es abierto, las imágenes $f(W)$, cuando $W$ recorre los
abiertos afines de $X$, forman un recubrimiento abierto de $Y$. Como $Y$
es cuasicompacto, existe un abierto cuasicompacto $T\subset X$
tal que $f(T)=Y$, y $f|T$ es cuasicompacto (I, 1.1.1). Con mayor razón,
para todo abierto cuasicompacto $T_\alpha\supset T$ de $X$,
\oldpage[IV]{360}
se tiene
$f(T_\alpha)=Y$ y $f|T_\alpha$ es cuasicompacto.

Aplicando el caso I a $f|T_\alpha$, obtenemos un divisor $D_\alpha$ sobre
$Y$ tal que
\[
 (f|T_\alpha)^*(D_\alpha)=D'|T_\alpha.
\]
Por transitividad de las imágenes recíprocas de divisores por morfismos planos
{\EGAReadLink{IV.21.4.4}{(21.4.4)}}, también se tiene
$(f|T)^*(D_\alpha)=D'|T$. Por tanto, para cualesquiera $\alpha,\beta$,
\[
 (f|T)^*(D_\alpha)=(f|T)^*(D_\beta),
\]
y {\EGAReadLink{IV.21.4.8}{(21.4.8)}} da $D_\alpha=D_\beta$. Si $D$ denota su valor
común, entonces $f^*(D)|T_\alpha=D'|T_\alpha$ para todo $\alpha$, y
por consiguiente $f^*(D)=D'$.
\end{proof}

\begin{corollary}[21.4.10]
\ifdefined\EGAIVSuppressErratumLabels\else\label{IV.21.4.10}\fi
Supongamos las hipótesis generales de {\EGAReadLink{IV.21.4.9}{(21.4.9)}} y, además, que
$Y$ es normal y que, para todo $y\in Y$ con
$\dim(\sh{O}_{Y,y})=1$, la fibra $X_y=f^{-1}(y)$ no tiene ningún
ciclo primo asociado inmerso. Un divisor $D'$ sobre $X$ es de la forma $f^*(D)$,
con $D\in\operatorname{Div}(Y)$, si y solo si se cumplen las dos
condiciones siguientes:
\begin{enumerate}[label=\textup{\arabic*\textdegree}]
\item para todo punto maximal $y$ de $Y$, la imagen recíproca de $D'$ en
$\operatorname{Div}(X_y)$ es cero;
\item para todo $y\in Y$ tal que $\dim(\sh{O}_{Y,y})=1$, existe un
entero $n_y$ tal que, para todo punto maximal $x$ de $X_y$, la
imagen recíproca de $D'$ en
$\operatorname{Div}(\operatorname{Spec}(\sh{O}_{X,x}))$ es igual a
$\operatorname{div}(t_y^{n_y})$, donde $t_y$ es la imagen por
$\sh{O}_{Y,y}\to\sh{O}_{X,x}$ de un uniformizante del anillo de
valoración discreta $\sh{O}_{Y,y}$.
\end{enumerate}
\end{corollary}

\begin{proof}
Aplicamos el criterio {\EGAReadLink{IV.21.4.9}{(21.4.9)}}. En un punto maximal $y$ de $Y$,
$\operatorname{Div}(Y_1)=0$ porque $\sh{O}_{Y,y}$ es artiniano;
la condición 1\textdegree{} es, por tanto, equivalente al criterio de
{\EGAReadLink{IV.21.4.9}{(21.4.9)}} en tal punto. Queda por demostrar que las condiciones
1\textdegree{} y 2\textdegree{} implican asimismo el criterio en un
punto $y\in Y$ con $\operatorname{prof}(\sh{O}_{Y,y})=1$. Como $Y$
es normal, estos son precisamente los puntos para los cuales
$\dim(\sh{O}_{Y,y})=1$, equivalentemente, aquellos para los cuales
$\sh{O}_{Y,y}$ es un anillo de valoración discreta {\EGAReadLink{IV.5.8.6}{(5.8.6)}}.

Podemos, pues, limitarnos a $Y=\operatorname{Spec}(\sh{O}_{Y,y})$. Si $t$
es un uniformizante de $\sh{O}_{Y,y}$, la aplicación
$n\mapsto\operatorname{div}(t^n)$ es un isomorfismo de $\mathbb Z$
sobre $\operatorname{Div}(Y)$. La condición de {\EGAReadLink{IV.21.4.9}{(21.4.9)}} dice que
$D'$ es de la forma $\operatorname{div}(t^n\circ f)$ para algún
$n\in\mathbb Z$, e implica evidentemente la condición 2\textdegree{}.
Recíprocamente, supongamos que se cumple 2\textdegree{}, así como 1\textdegree{}, y
escribamos $n=n_y$. Queda por demostrar
$D'=\operatorname{div}(t^n\circ f)$. Si $\eta$ es el punto genérico de
$Y$, los gérmenes de estos dos divisores en todo punto de $X_\eta$ son
cero por la condición 1\textdegree{}. Sus gérmenes en todo punto maximal
de $X_y$ son iguales por la condición 2\textdegree{}. En un punto no maximal
$x\in X_y$, se tiene
$\operatorname{prof}(\sh{O}_{X_y,x})\geq 1$, porque $X_y$ no tiene
ciclo primo asociado inmerso y, por tanto,
$\operatorname{prof}(\sh{O}_{X,x})\geq 2$ por {\EGAReadLink{IV.6.3.1}{(6.3.1)}}. La
igualdad buscada resulta de {\EGAReadLink{IV.21.1.8}{(21.1.8)}}.
\end{proof}

\begin{corollary}[21.4.11]
\ifdefined\EGAIVSuppressErratumLabels\else\label{IV.21.4.11}\fi
Supongamos las hipótesis generales de {\EGAReadLink{IV.21.4.9}{(21.4.9)}} y, además, que
$Y$ es normal y que, para todo $y\in Y$ con
$\dim(\sh{O}_{Y,y})=1$, la fibra $X_y$ es íntegra. Un divisor $D'$ es
de la forma $f^*(D)$, con $D\in\operatorname{Div}(Y)$, si y solo si,
para todo punto maximal $y$ de $Y$, la imagen recíproca de $D'$ en
$\operatorname{Div}(X_y)$ es cero.
\end{corollary}

\begin{proof}
Es suficiente comprobar que la condición 2\textdegree{} de {\EGAReadLink{IV.21.4.10}{(21.4.10)}} es entonces
automática. Si $\dim(\sh{O}_{Y,y})=1$ y $X_y$ es íntegra, entonces en
el único punto maximal $x$ de $X_y$ se tiene
$\dim(\sh{O}_{X,x})=1$ por {\EGAReadLink{IV.6.1.3}{(6.1.3)}}. Además,
$\sh{O}_{Y,y}$ es regular y $\sh{O}_{X_y,x}$ es un cuerpo, por lo que
$\sh{O}_{X,x}$ es regular {\EGAReadLink{IV.6.5.1}{(6.5.1)}} y, por tanto, un anillo de valoración discreta.
Si $t$ es un uniformizante de $\sh{O}_{Y,y}$, su imagen $t_y$ en
$\sh{O}_{X,x}$ es un uniformizante de $\sh{O}_{X,x}$
($0$, 17.3.3). La imagen recíproca de $D'$ en
$\operatorname{Div}(\operatorname{Spec}(\sh{O}_{X,x}))$ es, por tanto,
\oldpage[IV]{361}
de la forma $\operatorname{div}(t_y^n)$, lo que es exactamente la condición
2\textdegree{} de {\EGAReadLink{IV.21.4.10}{(21.4.10)}}, puesto que $X_y$ solo tiene un punto maximal.
\end{proof}

\begin{remark}[21.4.12]
\ifdefined\EGAIVSuppressErratumLabels\else\label{IV.21.4.12}\fi
Si en {\EGAReadLink{IV.21.4.11}{(21.4.11)}} se supone solamente que los $X_y$ son reducidos, hay que
añadir a la condición del enunciado que las imágenes recíprocas de $D'$
en cada uno de los grupos
$\operatorname{Div}(\operatorname{Spec}(\sh{O}_{X,x}))$, cuando $x$
recorre los puntos maximales de $X_y$, sean todas de la forma
$\operatorname{div}((t_x')^n)$, donde $t_x'$ es un uniformizante de
$\sh{O}_{X,x}$ y el entero $n$ es el mismo para todos estos puntos.
\end{remark}

\begin{corollary}[21.4.13]
\ifdefined\EGAIVSuppressErratumLabels\else\label{IV.21.4.13}\fi
Supongamos las hipótesis de {\EGAReadLink{IV.21.4.11}{(21.4.11)}}. Para todo
$\sh{O}_X$-módulo invertible $\sh{L}$ tal que, para todo punto maximal $\eta$ de
$Y$, la imagen recíproca
\[
 \sh{L}_\eta
 =\sh{L}\otimes_{\sh{O}_X}\sh{O}_{X_\eta}
\]
sea isomorfa a $\sh{O}_{X_\eta}$, existen un
$\sh{O}_Y$-módulo invertible $\sh{M}$ y un isomorfismo
$f^*(\sh{M})\simeq\sh{L}$.
\end{corollary}

\begin{proof}
Basta encontrar una sección meromorfa regular $s$ de $\sh{L}$ sobre
$X$ tal que el divisor $D'$ correspondiente a $(\sh{L},s)$ tenga
$D'_x=0$ en todo punto $x$ de una fibra $X_\eta$ sobre un punto maximal
$\eta$ de $Y$. Aplicando entonces {\EGAReadLink{IV.21.4.11}{(21.4.11)}} a $D'$, se obtiene el resultado por
{\EGAReadLink{IV.21.4.2}{(21.4.2)}}.

Los puntos de $\operatorname{Ass}(\sh{O}_X)$ están sobre los puntos
maximales de $Y$ {\EGAReadLink{IV.3.3.1}{(3.3.1)}}. Como $Y$ es reducido y localmente noetheriano,
existe un abierto denso $V$ de $Y$ que es la unión de una
familia disjunta de abiertos afines, cada uno de los cuales contiene exactamente un
punto maximal de $Y$; y $f^{-1}(V)$ es esquemáticamente denso en $X$
{\EGAReadLink{IV.5.10.2}{(5.10.2)}}. Podemos, por tanto, limitarnos al caso en que $Y$ es
afín e íntegro.

Para el único punto maximal $\eta$ de $Y$, la fibra $X_\eta$ es entonces
íntegra, de modo que $\operatorname{Ass}(\sh{O}_X)$ consta de un solo punto
$z$, el punto genérico de $X$. Sea $s_\eta$ la sección de
$\sh{L}_\eta$ sobre $X_\eta$ que corresponde, mediante un isomorfismo,
a la sección unidad de $\sh{O}_{X_\eta}$. Sea $U$ un entorno abierto
afín en $X$ de un punto $x\in X_\eta$, y sea $s$ una sección
de $\sh{L}$ sobre $U$ con el mismo germen en $x$ que $s_\eta$. Como $U$
es esquemáticamente denso en $X$ ({\EGAReadLink{IV.20.2.13}{20.2.13}}, (iv)), $s$ se identifica con una
sección meromorfa regular unívocamente determinada de $\sh{L}$ sobre $X$.
Resulta claramente de {\EGAReadLink{IV.20.2.11}{(20.2.11)}} y {\EGAReadLink{IV.20.3.5}{(20.3.5)}} que esta sección induce sobre
$X_\eta$ una sección meromorfa regular de $\sh{L}_\eta$ igual a
$s_\eta$ y, por tanto, tiene la propiedad requerida.
\end{proof}
